% !TEX root = ../Thesis.tex

\chapter{JHipster: Architettura e fondamenti teorici}

JHipster è una piattaforma open source progettata per la generazione di applicazioni web moderne full-stack. Il suo obiettivo è fornire una base pronta all’uso che integri tecnologie consolidate lato backend e frontend. Nel tempo, il progetto si è evoluto fino a diventare un ecosistema strutturato, mantenuto da una vasta comunità di sviluppatori \cite{JHipsterDocs}.

Attraverso una interfaccia a riga di comando, si risponde a una serie di domande che definiscono le caratteristiche principali dell’applicazione: tipologia di architettura (monolite o microservizi), sistema di autenticazione, database, strumenti di build e framework frontend, e alcuni strumenti addizionali. Si può in seguito, mediante la definizione di un modello di dominio tramite JDL (JHipster Domain Language), modellare entità, relazioni e vincoli in modo strutturato.

JHipster genera automaticamente una struttura completa di applicazione, comprensiva di backend REST, frontend SPA (Single Page Application), sistema di autenticazione, configurazione del database, test unitari e di integrazione, e strumenti di containerizzazione.

Le informazioni relative all'installazione e all'utilizzo di JHipster possono essere trovate sul sito ufficiale del progetto \url{www.jhipster.tech}. I progetti generati con JHipster supportano nativamente gli IDE più diffusi, come Eclipse utilizzato nel presente lavoro, e la loro documentazione spiega come utilizzare tutti i componenti integrati dal framework.

\section{Architettura backend: il ruolo di Spring Boot}

Il backend generato da JHipster utilizza Spring Boot, un framework dell'ecosistema Spring progettato per semplificare la configurazione e l’avvio di applicazioni Java. Spring Boot riduce la complessità di Spring tramite il principio di \textit{convention over configuration}, fornendo configurazioni predefinite che consentono di avviare rapidamente un’applicazione funzionante \cite{SpringBootDocs}.

L’architettura backend risultante è tipicamente stratificata in più livelli. Il \textit{domain layer} contiene le entità JPA che rappresentano il modello dati. Il \textit{repository layer}, basato su Spring Data JPA, fornisce l’accesso ai dati mediante interfacce che estendono \texttt{JpaRepository}. Il \textit{service layer} incapsula la logica di business, mentre il \textit{web layer} espone le API REST.

Questa separazione rende il codice più facilmente testabile, poiché consente di isolare la logica applicativa dal livello di persistenza e dal livello di presentazione. Dal punto di vista del TDD, la presenza di un service layer ben definito rappresenta un punto di ingresso naturale per la scrittura di test unitari.

\subsection{Profili di Spring Boot generati}
JHipster genera automaticamente diversi profili per Spring Boot. I profili permettono di definire configurazioni differenti dell’applicazione a seconda dell’ambiente di esecuzione.

In generale, i profili generati automaticamente da JHipster sono già configurati in modo adeguato per la maggior parte dei casi d’uso. Tuttavia, durante lo sviluppo di un’applicazione reale è spesso necessario personalizzare alcune configurazioni. Le modifiche più comuni riguardano principalmente:

\begin{itemize}
	\item la configurazione del database, ad esempio sostituendo il database embedded utilizzato in sviluppo con un database relazionale esterno;
	\item parametri di logging differenti tra ambiente di sviluppo e produzione;
	\item configurazioni di sicurezza e gestione delle sessioni;
	\item parametri di connessione a servizi esterni, come sistemi di messaggistica o storage.
\end{itemize}

JHipster genera un file di configurazione generale, \texttt{application.yml}, che contiene le proprietà comuni dell’applicazione, e due profili principali: \texttt{dev} e \texttt{prod}.

\begin{itemize}
	\item \textbf{application.yml}: file di configurazione base, comune a tutti gli ambienti. Contiene le proprietà generali dell’applicazione e viene integrato o sovrascritto dai profili specifici.
	
	\item \textbf{dev}: utilizzato durante lo sviluppo locale. Include configurazioni per logging più verboso, strumenti di sviluppo e un database locale.
	
	\item \textbf{prod}: destinato all’ambiente di produzione. Include configurazioni con maggiore attenzione a prestazioni, sicurezza e integrazione con servizi esterni.
\end{itemize}

Nel caso del profilo \texttt{dev}, JHipster attiva inoltre automaticamente il profilo accessorio \texttt{api-docs}, utile per l’esposizione della documentazione delle API durante lo sviluppo.

\section{Persistenza e gestione delle relazioni}

La persistenza dei dati è affidata a Spring Data JPA e Hibernate. Le entità vengono annotate secondo lo standard JPA e mappate su tabelle relazionali \cite{HibernateDocs}.

Uno degli aspetti più rilevanti è la gestione automatica delle relazioni tra entità, incluse relazioni molti-a-molti. Tramite JDL (o CLI interattiva) è possibile definire in modo dichiarativo entità, attributi, vincoli e relazioni (OneToOne, OneToMany, ManyToOne, ManyToMany), oltre a configurare DTO, service layer, paginazione e filtraggio. Il generatore produce automaticamente le annotazioni JPA appropriate e, nel caso delle relazioni molti-a-molti, crea la tabella di join necessaria.

Il Listing \ref{lst:jpa-entity} mostra un esempio di entità generata automaticamente dal framework.

\Needspace{12\baselineskip}
\captionsetup{type=listing}
\captionof{listing}{Esempio di entità JPA generata da JHipster}
\label{lst:jpa-entity}
\begin{minted}{java}
/**
* JHipster JDL model for myApp
*/
@Entity
@Table(name = "tag")
@Cache(usage = CacheConcurrencyStrategy.READ_WRITE)
@SuppressWarnings("common-java:DuplicatedBlocks")
public class Tag implements Serializable {
	
	private static final long serialVersionUID = 1L;
	
	@Id
	@GeneratedValue(strategy = GenerationType.IDENTITY)
	@Column(name = "id")
	private Long id;
	
	@NotNull
	@Size(min = 1, max = 255)
	@Column(name = "tag_name", length = 255, nullable = false)
	private String tagName;
	
	@Size(max = 20)
	@Column(name = "color", length = 20)
	private String color;
	
	@ManyToOne(optional = false)
	@NotNull
	private User owner;
	
	@ManyToMany(fetch = FetchType.LAZY, mappedBy = "tags")
	@Cache(usage = CacheConcurrencyStrategy.READ_WRITE)
	@JsonIgnoreProperties(value = { "attachments", "owner", "tags" }, allowSetters = true)
	private Set<Note> notes = new HashSet<>();
	
	...
}
\end{minted}

Le annotazioni come \texttt{@Entity}, \texttt{@Table}, \texttt{@Id} e \texttt{@GeneratedValue} permettono di definire la mappatura tra oggetti Java e tabelle del database. Questo consente a Spring Data JPA di generare automaticamente molte delle query necessarie per l’accesso ai dati. Le interfacce repository estendono infatti \texttt{JpaRepository}, che fornisce automaticamente implementazioni per le operazioni CRUD più comuni.

È inoltre possibile definire query semplicemente dichiarando metodi con nomi specifici all'interno dell'interfaccia repository. Spring Data JPA analizza la firma del metodo e genera automaticamente la query SQL corrispondente.

Il Listing \ref{lst:tag-repository} mostra un esempio di repository per l’entità \texttt{Tag}. 

\Needspace{12\baselineskip}
\captionsetup{type=listing}
\captionof{listing}{Esempio di repository Spring Data JPA generato da JHipster}
\label{lst:tag-repository}
\begin{minted}{java}
/**
* Spring Data JPA repository for the Tag entity.
*/
@Repository
public interface TagRepository extends JpaRepository<Tag, Long> {
	@Query("select tag from Tag tag where tag.owner.login = ?#{authentication.name}")
	List<Tag> findByOwnerIsCurrentUser();
	
	List<Tag> findDistinctAllByNotesId(@Param("noteId") Long noteid);
	
	Page<Tag> findAllByOwnerLogin(String login, Pageable pageable);
	
	Page<Tag> findDistinctAllByOwnerIdAndTagNameNotInAndTagNameStartingWithOrderByTagNameAsc(
	@Param("ownerId") Long ownerid,
	@Param("notetags") Collection<String> notetagname,
	@Param("initial") String initial,
	Pageable pageable
	);
	
	Page<Tag> findDistinctAllByOwnerIdAndTagNameStartingWithOrderByTagNameAsc(
	@Param("ownerId") Long ownerid,
	@Param("initial") String initial,
	Pageable pageable
	);
	
	@EntityGraph(attributePaths = "notes")
	Optional<Tag> findOneWithNotesByIdAndOwnerLogin(Long id, String login);
	
	@EntityGraph(attributePaths = "notes")
	Optional<Tag> findOneWithNotesById(Long id);
	
	Optional<Tag> findOneByIdAndOwnerLogin(@Param("id") Long id, @Param("login") String login);
	
	Optional<Tag> findOneByTagNameAndOwnerLogin(@Param("tagName") String tagName, @Param("login") String login);
	
	long countByIdInAndOwnerLogin(Collection<Long> ids, String login);
	
	boolean existsByTagNameAndOwnerLogin(@Param("tagName") String tagName, @Param("login") String login);
}
\end{minted}


\section{Architettura frontend: Angular}

JHipster supporta diversi framework frontend, tra cui: Angular, React, Vue. In questo progetto viene utilizzato Angular. L’applicazione generata da JHipster è una Single Page Application (SPA) che comunica con il backend esclusivamente tramite API REST.

Angular è un framework strutturato basato su TypeScript. Con Angular è possibile realizzare componenti, servizi e direttive che cooperano tra loro per costruire applicazioni client complesse. La comunicazione con il backend avviene tramite il modulo \texttt{HttpClient}, che gestisce richieste asincrone e serializzazione JSON \cite{AngularDocs}.

Il frontend generato include:

\begin{itemize}
	\item componenti CRUD per ciascuna entità;
	\item servizi dedicati alla comunicazione con le API REST;
	\item gestione del routing e delle guardie di autenticazione;
	\item integrazione con il sistema di sicurezza JWT.
\end{itemize}

\begin{figure}[H]
\dirtree{%
	.1 note/.
	.2 note.model.ts.
	.2 note.routes.ts.
	.2 note.test-samples.ts.
	.2 delete/.
	.3 note-delete-dialog.component.html.
	.3 note-delete-dialog.component.spec.ts.
	.3 note-delete-dialog.component.ts.
	.2 detail/.
	.3 note-detail.component.html.
	.3 note-detail.component.spec.ts.
	.3 note-detail.component.ts.
	.2 list/.
	.3 note.component.html.
	.3 note.component.spec.ts.
	.3 note.component.ts.
	.2 route/.
	.3 note-routing-resolve.service.spec.ts.
	.3 note-routing-resolve.service.ts.
	.2 service/.
	.3 note.service.spec.ts.
	.3 note.service.ts.
	.2 update/.
	.3 note-form.service.spec.ts.
	.3 note-form.service.ts.
	.3 note-update.component.html.
	.3 note-update.component.spec.ts.
	.3 note-update.component.ts.
}
\caption{Estratto dell’albero del frontend generato per una entità}
\label{fig:entity-ts-tree}
\end{figure}

La struttura evidenzia come JHipster organizzi automaticamente i componenti dell'interfaccia separando le responsabilità: la cartella \texttt{list} contiene i componenti per la visualizzazione delle entità, \texttt{detail} per la visualizzazione dei dettagli, \texttt{update} include i componenti e i servizi necessari alla creazione e modifica dei dati, mentre le cartelle \texttt{service} e \texttt{route} raccolgono rispettivamente la logica di comunicazione con il backend e quella di gestione del routing.

Dal punto di vista del testing, Angular supporta test unitari mediante strumenti come Jest, mentre per i test end-to-end JHipster integra Cypress.

\section{Autenticazione e autorizzazione}

La sicurezza dell’applicazione è gestita tramite Spring Security \cite{SpringSecurityDocs}. JHipster consente di scegliere tra vari metodi di autenticazione, i più semplici e supportati out of the box sono autenticazione basata sulla sessione e tramite JSON Web Tokens (JWT) \cite{RFC7519}.

La documentazione ufficiale di JHipster illustra inoltre come implementare e integrare gli altri metodi di autenticazione più complessi, e come renderli tutti più sicuri per l'ambiente di produzione.

Il sistema inoltre implementa ruoli predefiniti come \texttt{ROLE\_USER} e \texttt{ROLE\_ADMIN} per gli utenti. Le restrizioni di accesso vengono applicate sia a livello di endpoint REST sia a livello di frontend tramite guardie di routing.

\section{Supporto ai test}

Uno degli elementi distintivi di JHipster è la generazione automatica di test backend e frontend.

Per il backend vengono prodotti test basati su JUnit, Spring Test, MockMvc e Mockito. Ogni entità include test di integrazione per le operazioni CRUD, verifiche sulle validazioni e controlli sui codici di risposta HTTP.

Per il frontend Angular vengono generati test unitari per i componenti e servizi, e configurazioni per test end-to-end con Cypress.

Cypress è uno strumento moderno per test end-to-end che esegue il browser in modo controllato e consente di simulare interazioni. I test Cypress verificano il comportamento complessivo dell’applicazione, dall’interfaccia fino alla persistenza dei dati \cite{CypressDocs}.

Inoltre, al momento della creazione dell'applicazione, è possibile specificare l'utilizzo di test di performance tramite Gatling e test BDD tramite Cucumber.

Nel prossimo capitolo verranno analizzati più nel dettaglio i test generati automaticamente da JHipster e le modalità con cui è possibile estenderli per supportare un processo di sviluppo basato sul paradigma Test Driven Development (TDD).

\section{Containerizzazione e integrazione DevOps}

JHipster integra nativamente configurazioni Docker e file \texttt{docker-compose} per l’avvio coordinato di backend, frontend e database, con la possibilità di creazione tramite CLI. È inoltre possibile generare configurazioni per Kubernetes e file di pipeline per strumenti di integrazione GitHub Actions, GitLab CI o Jenkins.

La documentazione ufficiale di JHipster fornisce una spiegazione dettagliata su come realizzare il proprio Docker Compose e su come utilizzare i Docker Compose pregenerati per gli strumenti addizionali come Sonar.